\reportpage{\datacontinued}{
\taskheading{A1}{配置DHCP服务器与客户端}
\subheading{地址规划}
采用私有子网192.168.7.0/24，掩码255.255.255.0，网络地址为192.168.7.0，广播地址为192.168.7.255，可用主机范围为192.168.7.1～192.168.7.254。服务器配置静态地址192.168.7.2/24。

地址池名为LAB7，起始地址192.168.7.100，最大用户数10，因而覆盖192.168.7.100～192.168.7.109。该范围含10个可用地址，满足至少3台客户端的要求，避开服务器静态地址、网络地址和广播地址。

Server0、Switch0和前两台PC通过铜质直通线接入同一局域网。初始设备名PC0、PC1对应指导书中的前两台PC。默认网关和DNS保持未指定，实验通过同一子网内的通信验证分配结果。
\ev[100mm]{fig:a1-topology}{a1-topology}{A1初始拓扑及活动链路}
\ev[45mm]{fig:a1-server-ip}{a1-server-ip}{服务器的静态IPv4参数}
}
\reportpage{\datacontinued}{
\subheading{自动获取配置及连通性}
PC0、PC1分别获取192.168.7.100/24和192.168.7.101/24。对服务器的Ping均发送4包、接收4包。
\ev[42mm]{fig:a1-pool}{a1-pool}{DHCP开启后的地址池配置}
\ev[30mm]{fig:a1-ip0}{a1-pc0}{PC0的IPv4配置}
\ev[30mm]{fig:a1-ip1}{a1-pc1}{PC1的IPv4配置}
}
\reportpage{\datacontinued}{
\subheading{Ping验证}
\ev[60mm]{fig:a1-ping0}{a1-ping0}{PC0向服务器发送Ping}
\ev[60mm]{fig:a1-ping1}{a1-ping1}{PC1向服务器发送Ping}
}
\reportpage{\datacontinued}{
\subheading{问题分析}
\qa{6. PC1和PC2实际获得的地址是否符合地址池设计？说明依据。}{两台PC的地址互不相同，均落在192.168.7.100至192.168.7.109范围内，掩码均为255.255.255.0。界面同时列出默认池，按实际租约核对地址。}
\qa{7. 如果地址池只有一个可分配地址，两台PC都申请地址会怎样？为什么？}{在没有其他DHCP服务器、且唯一地址的租约仍有效的条件下，服务器只能把该地址有效租给其中一台客户端。另一台客户端无法从该池获得第二个互不冲突的租约，会继续请求或获取失败。部分操作系统在失败后会使用169.254.0.0/16范围的链路本地地址，但这不是从该DHCP地址池分配的地址；具体显示取决于客户端实现。本题未将地址池改成一个地址进行实测，以上为机制分析。}
}
